\section{Comparison for subdivided attachments}
\label{sec:attachments}

We now compare an ordinary pushout of Street nerves with the Street
nerve of an actual attachment.  All attaching functors in this section
are arbitrary.  In particular, they may identify old objects and cells.
The comparison will supply both the acyclicity of horn attachments and
the preservation of weak equivalences under boundary attachments.

We first treat a sieve $A\subset B$ whose inclusion into its upward
closure admits a right adjoint retraction.  Couniversality controls
the attachment over that upward closure; a cosieve cover then adds
the rest of $B$.  This decomposition follows the proof of
\cite[Proposition~4.3]{Thomason1980}; Theorem~\ref{thm:couniversal}
supplies its homotopy input in strict higher dimensions.  After proving
this general comparison, we identify the twice-subdivided boundaries
and horns as poset pairs with these properties.

\begin{proposition}[Attachment comparison]
\label{prop:attachment-comparison}
Let $A\subset B$ be a sieve of small posets.  Let $Z$ be its upward
closure, and suppose that $A\hookrightarrow Z$ admits a right adjoint
retraction.  For every functor $a:\U{A}\to X$ in $\nCat$, put
\[
 V_B(X)=X\amalg_{\U{A}}\U{B},\qquad
 P_B(X)=\Street X\amalg_{\Street\U{A}}\Street\U{B}.
\]
The canonical map
\begin{equation}
 \gamma_{B,X}:P_B(X)\longrightarrow\Street V_B(X)
 \label{eq:attachment-comparison}
\end{equation}
is a simplicial weak equivalence.  The map
$\Street\U{A}\to\Street\U{B}$ is a monomorphism.
The comparison is natural in every functor $f:X\to X'$ under
$\U{A}$.
\end{proposition}

\begin{proof}
First set
\[
 V_Z(X)=X\amalg_{\U{A}}\U{Z},\qquad
 P_Z(X)=\Street X\amalg_{\Street\U{A}}\Street\U{Z}.
\]
Theorem~\ref{thm:couniversal}, applied to $A\hookrightarrow Z$, says
that $X\to V_Z(X)$ belongs to $W_n$.  Applying it with old target
$\U{A}$ also says that $\U{A}\to\U{Z}$ belongs to $W_n$.
Since $A$ is a sieve in $Z$, Lemma~\ref{lem:poset-pieces} identifies
$\U{A}\to\U{Z}$ with an actual full inclusion.  Its Street nerve is
therefore a monomorphism and a weak equivalence, hence a trivial
cofibration of simplicial sets.  Its pushout
$\Street X\to P_Z(X)$ is a weak equivalence.  The canonical triangle
\[
 \begin{tikzcd}[column sep=large]
 \Street X \arrow[r] \arrow[dr] & P_Z(X) \arrow[d,"\gamma_{Z,X}"]\\
 & \Street V_Z(X)
 \end{tikzcd}
\]
then shows by two-out-of-three that $\gamma_{Z,X}$ is a weak
equivalence.

To pass from $Z$ to $B$, write
\[
 C=B\setminus A,\qquad T=Z\setminus A=Z\cap C.
\]
Both $Z$ and $C$ are cosieves in $B$, and they cover $B$: an element
outside $C$ belongs to $A\subset Z$.  Their intersection is $T$.
Lemma~\ref{lem:poset-pieces} therefore gives cosieves
$\U{Z},\U{C}\subset\U{B}$ with intersection $\U{T}$.
Apply Corollary~\ref{cor:attachment-covers}(ii) with $P=B$,
$T_1=Z$, and $T_2=C$.  Its remaining hypotheses hold because
$A\subset Z$ and $A\cap C=\varnothing$.  The canonical maps identify
\begin{equation}
 V_Z(X),\quad \U{C}\ \subset\ V_B(X),
 \qquad V_Z(X)\cap\U{C}=\U{T},
 \label{eq:attachment-fibers}
\end{equation}
as actual full cosieves and their full intersection.  These two
cosieves cover all objects of $V_B(X)$.  The corollary applies to the
given arbitrary functor $a$, including its identifications of old
objects and cells.

Apply Lemma~\ref{lem:nerve-cover} to the two cosieve covers just
obtained.  It gives ordinary simplicial pushouts
\begin{align}
 \Street\U{B}
   &=\Street\U{Z}\amalg_{\Street\U{T}}\Street\U{C},
   \label{eq:attachment-source-cover}\\
 \Street V_B(X)
   &=\Street V_Z(X)\amalg_{\Street\U{T}}\Street\U{C}.
   \label{eq:attachment-target-cover}
\end{align}
All maps from the intersections to the pieces in these covers are
monomorphisms.  Reassociating the pushout in the definition of
$P_B(X)$ using Equation~\eqref{eq:attachment-source-cover} gives
\[
 P_B(X)=P_Z(X)\amalg_{\Street\U{T}}\Street\U{C}.
\]
Under these canonical identifications, Equation~\eqref{eq:attachment-comparison}
is exactly the map induced by $\gamma_{Z,X}$ and the identity of
$\Street\U{C}$.  In particular the square
\[
 \begin{tikzcd}[column sep=large]
 P_Z(X) \arrow[r,hook] \arrow[d,"\gamma_{Z,X}"'] &
 P_B(X) \arrow[d,"\gamma_{B,X}"]\\
 \Street V_Z(X) \arrow[r,hook] & \Street V_B(X)
 \end{tikzcd}
\]
is a pushout: its lower right corner is also
$\Street V_Z(X)\amalg_{P_Z(X)}P_B(X)$.
The top horizontal map is the cobase change of the specific
monomorphism
\begin{equation}
 \Street\U{T}\hookrightarrow\Street\U{C}.
 \label{eq:comparison-monomorphism}
\end{equation}
Left properness of simplicial sets, as recorded in
Lemma~\ref{lem:simplicial-gluing}, now makes $\gamma_{B,X}$ a weak
equivalence.  Neither the map $\Street\U{A}\to\Street X$ nor the
map $\Street\U{T}\to P_Z(X)$ needs to be a monomorphism.

The remaining assertion that $\Street\U{A}\to\Street\U{B}$ is
monic follows directly from the sieve identification in
Lemma~\ref{lem:poset-pieces}.  Finally, $\gamma_{B,X}$ is defined
by the two canonical maps from $\Street X$ and $\Street\U{B}$.
For $f:X\to X'$ under $\U{A}$ these maps commute with $\Street f$
and with the induced pushout functor $V_B(X)\to V_B(X')$.
The pushout universal property therefore gives the strict naturality
square
\begin{equation}
 \begin{tikzcd}[column sep=large]
 P_B(X) \arrow[r,"\gamma_{B,X}"] \arrow[d] &
 \Street V_B(X) \arrow[d]\\
 P_B(X') \arrow[r,"\gamma_{B,X'}"'] & \Street V_B(X').
 \end{tikzcd}
 \label{eq:attachment-naturality}
\end{equation}
The comparison is thus canonical and natural, while the local
contractions used in Theorem~\ref{thm:couniversal} establish that it
is a weak equivalence.
\end{proof}

To apply the comparison to the generating maps, we now describe their
second subdivisions.  We use the classical face-chain description and
intersection retraction of \cite[Proposition~4.2]{Thomason1980}; see also
\cite[Lemmas~4.8 and~4.10]{AraMaltsiniotis2014}. The proof below records
the precise poset pair for each boundary and horn, including the empty case.

\begin{lemma}[The subdivision posets]
\label{lem:subdivision-posets}
Let $d:K\hookrightarrow\Delta[m]$ be a boundary inclusion, with $m\geq0$,
or a horn inclusion $\Lambda^k[m]\hookrightarrow\Delta[m]$, with
$m\geq1$ and $0\leq k\leq m$.  There are finite posets $A\subset B$
and a canonical identification
\[
 \Sd^2 d=\bigl(\OrdN A\hookrightarrow\OrdN B\bigr).
\]
Here $A$ is a sieve.  Its upward closure $Z$ in $B$ is a cosieve, and
$A\hookrightarrow Z$ has a right adjoint retraction.
\end{lemma}

\begin{proof}
Put $M=\{0,\ldots,m\}$ and let $P$ be the poset of its nonempty
subsets, ordered by inclusion.  The subposet $P_K$ of faces belonging
to $K$ is
\[
 P_K=
 \begin{cases}
  \{F:\varnothing\ne F\subsetneq M\},&K=\partial\Delta[m],\\[2pt]
  \{F:\varnothing\ne F\subseteq M,\ M\setminus\{k\}\nsubseteq F\},
     &K=\Lambda^k[m].
 \end{cases}
\]
Indeed, a nonempty face belongs to the horn precisely when it is
contained in $M\setminus\{j\}$ for some $j\ne k$.  Thus in the horn
case $P_K$ is obtained from $P$ by removing $M$ and
$M\setminus\{k\}$.  In both cases $P_K$ is downward closed in $P$.

The barycentric subdivision of an ordered simplicial complex has a
vertex for every nonempty face and a nondegenerate simplex for every
strict chain of such faces.  Consequently
\[
 \Sd\Delta[m]=\OrdN P,\qquad \Sd K=\OrdN P_K.
\]
One can see these identifications directly on a simplex: the vertices
of its subdivision are its nonempty subsets, and its simplices are
chains of those subsets.  The identifications agree on every face, so
they glue over the faces of $K$.  The nondegenerate simplices of
$\OrdN P$ are themselves the nonempty finite chains in $P$, with a
face obtained by deleting entries.  Applying the same description a
second time gives
\[
 \begin{split}
 B&=\{\sigma:\sigma\text{ is a nonempty chain in }P\},\\
 A&=\{\sigma\in B:\sigma\subseteq P_K\},
 \end{split}
 \qquad \sigma\leq\tau\ \Longleftrightarrow\ \sigma\subseteq\tau.
\]
This proves the claimed identifications as simplicial sets, including
degeneracies, and identifies their inclusion.

Every nonempty subchain of a chain in $A$ again belongs to $A$, so $A$
is a sieve in $B$.  Its upward closure is exactly
\[
 Z=\{\sigma\in B:\sigma\cap P_K\ne\varnothing\}.
\]
For the nontrivial implication in this equality, a face
$F\in\sigma\cap P_K$ gives $\{F\}\in A$ below $\sigma$.
Define
\[
 r:Z\longrightarrow A,\qquad r(\sigma)=\sigma\cap P_K.
\]
This is a nonempty chain, is monotone under inclusion, and fixes $A$.
For $\alpha\in A$ and $\sigma\in Z$ one has
\[
 \alpha\subseteq\sigma
 \quad\Longleftrightarrow\quad
 \alpha\subseteq\sigma\cap P_K=r(\sigma).
\]
This is precisely the right adjunction to $A\hookrightarrow Z$.

When $m=0$ and $K=\partial\Delta[0]$, the poset $B$ is a singleton
and $A=Z=\varnothing$; the asserted adjunction is the identity of the
empty poset.  For a one-dimensional horn, $P_K=\{\{k\}\}$, so $A$
is a singleton and $Z$ consists of the chains containing the face
$\{k\}$.  These cases are included in the formulas above.  There are
no zero-dimensional horn generators.
\end{proof}

\begin{remark}
For the zero-dimensional boundary, $A=Z=T=\varnothing$ and $B=C$
is a singleton.  Thus $V_Z(X)=X$ and $\gamma_{Z,X}=\id$; the two
cosieves of $V_B(X)=X\amalg\U{B}$ are disjoint.  The comparison in
Proposition~\ref{prop:attachment-comparison} is then the isomorphism
\[
 \Street X\amalg\Delta[0]\ \cong\ \Street(X\amalg\U{B}).
\]
\end{remark}

\begin{corollary}[Horn attachments are acyclic]
\label{cor:horn-acyclic}
For every horn inclusion $d\in J$, every pushout of $L_n d$ belongs
to $W_n$.
\end{corollary}

\begin{proof}
Use the pair $A\subset B$ from Lemma~\ref{lem:subdivision-posets},
so that $L_n d$ is $\U{A}\to\U{B}$.
The inclusion $\OrdN A\to\OrdN B$ is a weak equivalence: the horn
inclusion is a simplicial weak equivalence, and subdivision preserves
weak equivalences.  The latter follows by naturality and
two-out-of-three from the objectwise weak equivalence
$\Sd\to\id$ of \cite[Theorem~4.5]{AraMaltsiniotis2014}.
Naturality of the adjunction unit gives a square
\[
 \begin{tikzcd}[column sep=large]
 \OrdN A \arrow[r] \arrow[d] & \OrdN B \arrow[d]\\
 \Street\U{A} \arrow[r] & \Street\U{B}.
 \end{tikzcd}
\]
The vertical maps are weak equivalences by
Proposition~\ref{prop:poset-unit}.  Two-out-of-three makes the lower
map a weak equivalence, and it is monic by
Proposition~\ref{prop:attachment-comparison}.  Thus it is a
simplicial trivial cofibration.  For any attaching functor
$\U{A}\to X$, its pushout $\Street X\to P_B(X)$ is a weak
equivalence.  The composite
\[
 \Street X\longrightarrow P_B(X)
   \xrightarrow{\gamma_{B,X}}\Street V_B(X)
\]
is the nerve of the attachment $X\to V_B(X)$, and both factors are
weak equivalences.  This is exactly the assertion that the
attachment belongs to $W_n$.
\end{proof}

\begin{corollary}[Invariance under attachment]
\label{cor:attachment-invariance}
Let $d\in I\cup J$, write $L_n d$ as $\U{A}\to\U{B}$ using
Lemma~\ref{lem:subdivision-posets}, and let $f:X\to X'$ belong to
$W_n$.  For every $a:\U{A}\to X$, the induced functor
\[
 X\amalg_{\U{A}}\U{B}
 \longrightarrow X'\amalg_{\U{A}}\U{B},
\]
where the second attaching functor is $fa$, belongs to $W_n$.
\end{corollary}

\begin{proof}
The map $\Street X\to P_B(X)$ is a monomorphism, being the
cobase change of $\Street\U{A}\hookrightarrow\Street\U{B}$.
There is a canonical pushout identification
\[
 P_B(X')=\Street X'\amalg_{\Street X}P_B(X).
\]
Left properness of simplicial sets therefore makes
$P_B(X)\to P_B(X')$ a weak equivalence, since $\Street f$ is one.
In the naturality square of Equation~\eqref{eq:attachment-naturality}, both horizontal
maps are weak equivalences by
Proposition~\ref{prop:attachment-comparison}.  Two-out-of-three
implies that its right vertical map is a weak equivalence.  This is
the nerve of the asserted functor.
\end{proof}
